\documentclass[a4paper, serif, 10pt]{moderncv}

%% moderncv
% style and color
\moderncvstyle{banking}
\moderncvcolor{black}

%% page layout/margins
\usepackage[top=2.5cm, bottom=2.5cm, left=1.5cm, right=1.5cm]{geometry}

\nopagenumbers{}

%% Babel config for English language
\usepackage[english]{babel}

%% fontspec
\usepackage{fontspec}

\IfFontExistsTF{Linux Libertine}{
  \setmainfont[Ligatures={TeX, Common}, Numbers=OldStyle]{Linux Libertine}
}{}
\IfFontExistsTF{Linux Libertine O}{
  \setmainfont[Ligatures={TeX, Common}, Numbers=OldStyle]{Linux Libertine O}
}{}

%% CTeX
% CJK support, used to show CJK characters in the resume
%
% - fontset=none: disable builtin fontset but instead set the CJK font manually
% - heading=false: disable ctex heading
% - punct=kaiming: use kaiming punctuations styles for CJK
% - scheme=plain: use plain scheme, do not override \normalsize font size
% - space=auto: space settings for CJK characters
%
% ref:
% - http://ctan.mirrorcatalogs.com/language/chinese/ctex/ctex.pdf
\usepackage[UTF8, heading=false, punct=kaiming, scheme=plain, space=auto]{ctex}

\IfFontExistsTF{Noto Serif CJK SC}{
  \setCJKmainfont{Noto Serif CJK SC}
}{}
\IfFontExistsTF{Noto Sans CJK SC}{
  \setCJKsansfont{Noto Sans CJK SC}
}{}

%% line spacing
\usepackage{setspace}
\setstretch{1.125}

%% URL styling - use normal text instead of monospace
\urlstyle{same}

% Auto-underline all links
% Defer until \begin{document} so hyperref is already loaded
\AtBeginDocument{%
  \let\oldhref\href
  \renewcommand{\href}[2]{\oldhref{#1}{\underline{#2}}}%
}

%% Patch \httplink and \httpslink to handle full URLs
% The original moderncv macros always prepend http:// or https:// to URLs.
% This causes issues when the URL already has a protocol (e.g., https://example.com)
% resulting in malformed URLs like https://https://example.com.
% This patch checks if the URL already contains "://" and uses it directly if so.
% Note: We use \str_set:Nx (x-expansion) to expand macros like \@homepage first.
\ExplSyntaxOn
\str_new:N \g_yamlresume_url_str

\RenewDocumentCommand{\httpslink}{O{}m}{
  \str_set:Nx \g_yamlresume_url_str {#2}
  \str_if_in:NnTF \g_yamlresume_url_str {://}
    { \url{#2} }
    { \url{https://#2} }
}

\RenewDocumentCommand{\httplink}{O{}m}{
  \str_set:Nx \g_yamlresume_url_str {#2}
  \str_if_in:NnTF \g_yamlresume_url_str {://}
    { \url{#2} }
    { \url{http://#2} }
}
\ExplSyntaxOff

\name{Elena Marsh}{}
\title{UX Designer creating inclusive, human-centered digital experiences}
\phone[mobile]{(415) 555-0139}
\email{elena.marsh@example.com}
\homepage{https://elenamarsh.design}

\address{548 Market St, Suite 300, San Francisco, California, United States, 94104}{}{}

\extrainfo{{\small \faLinkedin}\ \href{https://www.linkedin.com/in/elenamarshux}{@elenamarshux} {} {} {} • {} {} {}
{\small \faDribbble}\ \href{https://dribbble.com/elenamarsh}{@elenamarsh} {} {} {} • {} {} {}
{\small \faMedium}\ \href{https://medium.com/@elenamarsh.ux}{@elenamarsh.ux}}

\begin{document}

\maketitle

\section{Basics}

\cvline{}{\begin{itemize}
\item UX designer with 7+ years of experience designing web and mobile products across health tech, fintech, and analytics
\item Human-centered design advocate specializing in user research, prototyping, and design systems
\item Collaborative team player skilled at partnering with product managers, engineers, and stakeholders to ship intuitive solutions
\item Passionate about accessibility and inclusive design, ensuring every user can accomplish their goals with ease
\end{itemize}}
\section{Education}

\cventry{Sep 2013–May 2017}
        {Bachelor, Interaction Design, Score: 3.8}
        {California College of the Arts}
        {\url{https://www.cca.edu/}}
        {}
        {\begin{itemize}
\item Developed a strong foundation in user-centered design principles and iterative design processes
\item Completed hands-on projects in mobile app design, responsive web design, and interactive installations
\item Recognized for outstanding work in the senior capstone project, a public transit navigation tool for elderly users
\end{itemize}
\textbf{Courses}: Human-Computer Interaction, Visual Communication, Design Research Methods, Prototyping for Digital Products, Psychology of User Experience, Design Ethics and Accessibility}
\section{Work}

\cventry{Mar 2021–Present}
        {Senior Product Designer}
        {Nimbus Analytics}
        {\url{https://www.nimbusanalytics.com/}}
        {}
        {\begin{itemize}
\item Led end-to-end design for a data visualization platform used by 50,000+ analysts, revamping the dashboard experience and reducing onboarding time by 35\%
\item Built and maintained the company’s first design system, including 120+ accessible components and tokens, resulting in a 40\% faster handoff to engineering
\item Drove user research initiatives, conducting 30+ contextual interviews and usability tests to shape the product roadmap
\item Partnered with product and engineering leadership to define design strategy and quarterly OKRs for the analytics suite
\item Mentored two junior designers, facilitating design critiques and skill-building workshops
\end{itemize}
\textbf{Keywords}: Product Strategy, Design Systems, User Research, Data Visualization, Figma, Mentorship}

\cventry{Jun 2017–Feb 2021}
        {Product Designer}
        {Brightline Health}
        {\url{https://www.brightlinehealth.com/}}
        {}
        {\begin{itemize}
\item Designed patient-facing mobile app and web portal for a telehealth startup, improving task completion rate by 28\%
\item Collaborated with clinical and engineering teams to translate complex healthcare workflows into simple, accessible experiences
\item Planned and moderated remote usability testing sessions, synthesizing findings into actionable design recommendations
\item Created high-fidelity prototypes in Figma and ran rapid iteration cycles to validate concepts with real patients
\item Advocated for accessibility standards, partnering with engineers to meet WCAG 2.1 AA compliance across all products
\end{itemize}
\textbf{Keywords}: Healthcare, Mobile App Design, Accessibility, Prototyping, Usability Testing, User Flows}
\section{Languages}

\cvline{English}{Native or Bilingual Proficiency \hfill \textbf{Keywords}: Native speaker, Writing and documentation}
\cvline{Spanish}{Limited Working Proficiency \hfill \textbf{Keywords}: Conversational, Basic reading comprehension}
\section{Skills}

\cvline{User Research}{Advanced \hfill \textbf{Keywords}: Interviews, Surveys, Personas, Journey Mapping, Usability Testing}
\cvline{Prototyping}{Expert \hfill \textbf{Keywords}: Figma, Sketch, Adobe XD, Framer, Interactive Prototypes}
\cvline{Visual Design}{Advanced \hfill \textbf{Keywords}: Typography, Color Theory, Layout, Design Systems, Accessibility (WCAG)}
\cvline{UX Strategy}{Intermediate \hfill \textbf{Keywords}: Product Roadmapping, Stakeholder Management, Design Sprints, Metrics, A/B Testing}
\section{Awards}

\cventry{Oct 2022}
        {Interaction Design Association}
        {UX Design Award — Best Mobile App}
        {}
        {}
        {Recognized for designing a mindful spending app that helps users understand their financial habits through visual data and gentle behavioral nudges.}
\section{Certificates}

\cventry{Jun 2022}
        {Nielsen Norman Group}
        {NN/g UX Certification}
        {\url{https://www.nngroup.com/}}
        {}
        {}

\cventry{Aug 2021}
        {Google}
        {Google UX Design Professional Certificate}
        {\url{https://www.coursera.org/professional-certificates/google-ux-design}}
        {}
        {}
\section{Publications}

\cventry{Jan 2021}
        {Designing for Trust in Health Tech}
        {UX Collective}
        {\url{https://uxdesign.cc/designing-for-trust-in-health-tech}}
        {}
        {Explores how transparency, clear language, and coherent visual hierarchy can build trust in digital health products, sharing lessons learned from patient-facing design research.}
\section{Projects}

\cventry{Jan 2020–May 2020}
        {A personal finance app concept that helps users build mindful spending habits through visual spending insights and gentle, non-judgmental nudges.}
        {Mindful Spending App}
        {\url{https://www.behance.net/gallery/mindful-spending-app}}
        {}
        {\begin{itemize}
\item Conducted diary studies and interviews with 12 participants to understand emotional relationships with money
\item Designed low-fidelity paper prototypes and iterated through high-fidelity Figma mockups based on feedback
\item Developed an interactive prototype in Framer to test navigation and core-spending reflection features
\item Awarded an honorable mention in the IxDA student design showcase
\end{itemize}
\textbf{Keywords}: Mobile App, Fintech, User Research, Prototyping, Accessibility}
\section{Interests}

\cvline{Photography}{Street photography, Portraits, Film photography}
\cvline{Accessibility}{Inclusive design, Assistive technology, Web accessibility}
\cvline{Hiking}{Bay Area trails, National parks, Nature photography}
\section{Volunteer}

\cventry{Sep 2018–Dec 2022}
        {Design Mentor}
        {Design4Good SF}
        {\url{https://www.design4goodsf.org/}}
        {}
        {\begin{itemize}
\item Mentored early-career designers from underrepresented backgrounds in portfolio development and interview preparation
\item Facilitated monthly critique sessions to help mentees refine their design thinking and presentation skills
\item Partnered with local nonprofits on pro-bono design projects, including a website refresh for a food bank and a donation flow redesign for a homeless shelter
\end{itemize}}

\end{document}